We model \emph{Dancing with the Stars} as a replayable inverse problem: judges' scores and eliminations are observed,
while audience votes are latent. Task~A infers weekly fan vote shares consistent with observed exits and reports
identifiability/uncertainty; Tasks~B--C use these inferred votes to replay eliminations across seasons to quantify how
aggregation mechanics (rank vs.\ percent) and the incumbent Bottom-2 + judges' save safeguard change outcomes. Any
completion needed when a counterfactual keeps an observed-exit contestant active is handled conservatively for replay
closure only (not for forecasting).

Building on the inferred fan signal and mechanism interpretation (Task~D1), Task~D2 proposes \emph{Professional Bottom-2} (PB2):
judges first define a small performance-grounded risk set, and fan influence is applied only within that gate via a
transparent weighted blend. PB2 preserves audience agency while reducing extreme reversals, and we validate it through
replay-based stress tests (stability across rules and seasons, controversial-case replays, and a unified bias check).

Overall, the pipeline provides an auditable chain from cleaned data to interpretable mechanisms to actionable rule
recommendations, while explicitly flagging weeks/seasons where non-identifiability limits confidence. Future work can
incorporate additional popularity signals and evaluate alternative safeguards within the same replay-and-stress-test
framework.
